\FloatBarrier
\newpage
\section{条件概率}
\label{sec:prob-conditional}

已经知道一件事发生，相当于获得了关于试验结果的新信息。概率计算需要在这部分结果中重新分配权重。条件可能来自试验的先后过程，也可能来自对同一次结果的进一步观察。

\subsection{条件下的模型}

\begin{definition}
若\(P(B)>0\)，称
\begin{equation}\label{eq:prob-conditional}
 P(A\mid B)=\frac{P(A\cap B)}{P(B)}
\end{equation}
为已知\(B\)发生时\(A\)的{\heiti 条件概率}。
\end{definition}

在有限模型中，保留\(B\)内的样本点，将其原权重都除以\(P(B)\)，总权重便重新成为1。\(A\)在这个模型中的部分是\(A\cap B\)，因此得到上式。若原样本点等可能，还可写为\(P(A\mid B)=|A\cap B|/|B|\)，即在\(B\)内重新计数。

例如公平抛两次硬币，已知至少一次正面，剩下\(HH,HT,TH\)三个等可能结果，其中两次正面的概率为\(1/3\)。若已知第一次正面，剩下的是\(HH,HT\)，条件概率则为\(1/2\)。不同信息对应不同条件事件。

\begin{example}[2023年全国甲卷理科，题面整理]
某地中学生中，60\%爱好滑冰，50\%爱好滑雪，70\%至少爱好其中一项。随机调查一名学生，已知他爱好滑雪，求他也爱好滑冰的概率。
\end{example}

\subsection{乘法公式}

把定义移项，即得
\begin{equation}\label{eq:prob-multiplication}
 P(A\cap B)=P(B)P(A\mid B).
\end{equation}
多阶段事件可以连续使用此式。例如
\[
 P(A\cap B\cap C)=P(A)P(B\mid A)P(C\mid A\cap B),
\]
其中各条件事件均有正概率。树状图的每条完整路径，把沿途的条件概率相乘；不同终点路径互斥时，再把路径概率相加。

袋中有3个红球、2个白球，不放回抽两球，所得两球异色的概率为
\[
 \frac35\frac24+\frac25\frac34=\frac35.
\]
第二次抽取的条件概率由第一次实际拿走的颜色决定。放回抽取时，袋中构成恢复，路径上的相应分母和分子便随之改变。

\begin{example}[2022年天津卷，题面整理]
从含4张A的52张扑克牌中依次随机取两张，不放回。求两张都是A的概率，以及第一张是A时第二张也是A的条件概率。
\end{example}

\subsection{全概率公式}

\begin{theorem}[全概率公式]
设\(B_1,\ldots,B_m\)两两互斥，其并为\(\Omega\)，且\(P(B_i)>0\)，则
\begin{equation}\label{eq:prob-total}
 P(A)=\sum_{i=1}^mP(B_i)P(A\mid B_i).
\end{equation}
\end{theorem}

\begin{proof}
\(A\)是互斥事件\(A\cap B_i\)的并，先用可加性，再用乘法公式即可。
\end{proof}

下面固定一个抽检模型供本节使用。产品来自甲、乙两条生产线，来源比例分别为\(0.4,0.6\)，不合格率分别为\(0.02,0.05\)。从混合产品中随机抽一件，令\(D\)表示不合格，则
\[
 P(D)=0.4\cdot0.02+0.6\cdot0.05=0.038.
\]
总体不合格率是各来源不合格率按来源比例的加权平均。图中路径同时保留了来源与检验结果。

\begin{figure}[!ht]
\centering
\includegraphics[width=.58\linewidth]{probability-tree}
\caption{按来源分类，再记录检验结果；每条路径的概率由两个条件因子相乘得到。}
\end{figure}

\subsection{贝叶斯公式}

\begin{theorem}[贝叶斯公式]
在全概率公式的条件下，若\(P(A)>0\)，则
\begin{equation}\label{eq:prob-bayes}
 P(B_j\mid A)=
 \frac{P(B_j)P(A\mid B_j)}{\sum_iP(B_i)P(A\mid B_i)}.
\end{equation}
\end{theorem}

\begin{proof}
在条件概率定义中代入乘法公式与全概率公式即得。
\end{proof}

在上述抽检模型中，已知产品不合格，它来自甲线的概率为
\[
 P(\text{甲}\mid D)=\frac{0.4\cdot0.02}{0.038}=\frac4{19}.
\]
抽检以前，甲线占全部产品的\(0.4\)；知道不合格后，应比较各线贡献的不合格产品数，比例于是改变。分母负责把这些贡献重新归一化。

把上述比例换成每万件产品的构成，可以更直接地看到所取的分母：
\[
\begin{array}{c|rr|r}
 &\text{不合格}&\text{合格}&\text{合计}\\\hline
 \text{甲线}&80&3920&4000\\
 \text{乙线}&300&5700&6000\\\hline
 \text{合计}&380&9620&10000
\end{array}
\]
表中数值按模型比例列出。甲线有80件不合格产品，已知来自甲线时分母取4000，已知不合格时则取380。分子对应同一交事件，分母随条件改变。

\begin{example}[据2025年北京卷改编]
甲、乙两校各随机调查100名学生，分别有80人、75人答对某四选一选择题。用频率估计答对概率。甲校学生掌握知识点时必答对，乙校学生掌握时答对概率为\(0.85\)；两校未掌握者均等可能地猜选四个选项。分别估计两校学生掌握知识点的比例，并比较大小。
\end{example}

\begin{exercise}
\mbox{}\par\nobreak

（1）某盒中有4个红球、3个白球，依次不放回抽两个。已知第二个是白球，求第一个是红球的概率。

（2）在本节生产线模型中，已知抽到的产品合格，求它来自乙线的概率。

（3）盒甲有2个红球、1个白球，盒乙有1个红球、3个白球。等可能地选一盒，再从中随机取一球。已知取出红球，求所选为盒甲的概率。若改为从全部7个球中等可能抽一个，条件概率如何改变？

（4）设\(P(B)=0.3\)、\(P(A\mid B)=0.8\)、\(P(A\mid B^c)=0.2\)。求\(P(A)\)、\(P(B\mid A)\)及\(P(B\mid A^c)\)。

（5）（据2024年天津卷改编）从包含“整地做畦”和“田间灌溉”的5个劳动技能项目中等可能选3项。求选中前一项的概率，以及已选中前一项时也选中后一项的概率。
\end{exercise}

\FloatBarrier
\newpage
\subsection{信息与分割*}

将\(\Omega\)分成两两不交且非空的\(B_1,\ldots,B_m\)，就得到一个{\heiti 分割}。若只记录样本点属于哪个\(B_i\)，同一组内的结果在当前信息下便无法区分。对\(\Omega\)的每个样本点都单独记录，给出最细的分割；只记录必然事件，则没有获得新的区分。

例如，抛两次硬币时，“正面总次数”给出分割
\[
 \{HH\},\quad\{HT,TH\},\quad\{TT\};
\]
“第一次的面”给出\(\{HH,HT\},\{TH,TT\}\)。这两种记录各保留一部分信息，前者能确定总次数，后者能确定先后位置。若同时记录两者，就能区分四个结果。

若分割\(\mathcal Q\)中的每一组都包含于分割\(\mathcal P\)的某一组，称\(\mathcal Q\)比\(\mathcal P\){\heiti 更细}。从粗记录获得更细记录，原来的条件概率可以继续按组分解。例如\(B=C_1\cup\cdots\cup C_r\)是互斥并，则
\[
 P(A\mid B)=\sum_j P(C_j\mid B)P(A\mid C_j).
\]
这是在条件模型\(B\)内再次使用全概率公式。后面按信息分组求条件平均，也将遵循相同的结构。

\begin{exercise}
掷两枚骰子，设36个有序点数对等可能。分别考虑记录“点数和”“第一枚点数”“两个点数的有序对”所形成的分割。比较其粗细关系，并求已知点数和为8时第一枚大于3的概率。
\end{exercise}
